\section{A boundary with a quasi-projective neighborhood}
We first prove the projective case and then keep the entire boundary
fixed during the passage to a projective model.
The projective input follows from the curve theorem of
Drinfeld--Esnault--Kindler \cite[Lemma C.2, p.~19]{Drinfeld}
\cite[Proposition 6.2, pp.~8--9]{EsnaultKindler}: enlarge the boundary
to a divisor containing the singular locus and its components of
codimension at least two. A cover of the original open is unramified
at each added generic divisor. The larger-boundary quotient therefore
surjects onto the original one. The proof below writes out the fixed-curve
argument in this closed-boundary formulation.

\begin{lemma}\label{lem:projective}
For a normal connected projective variety $X$ over an algebraically
closed field and a proper closed subset $Y$, the group $Q(X,Y)$ is
topologically finitely generated.
\end{lemma}
\begin{proof}
The zero-dimensional case is immediate. For a curve this is the
tame-curve theorem \cite[XIII, Corollary 2.12,
pp.~392--395]{SGA1}. Suppose $d=\dim X\ge2$, and embed
$X$ in $\mathbf P^r$. The singular locus of $X$, the components
of $Y$ of codimension at least two, and the singular and intersection
loci of its reduced divisorial components form a closed subset
$T\subset X$ of codimension at least two. Ordinary Bertini and
the dimension count for $T$ give a linear subspace $L$ of
codimension $d-1$ such that $C=X\cap L$ is smooth, avoids $T$,
and meets the boundary divisors transversally. These are finitely
many nonempty open conditions on the parameter space over the
algebraically closed field. Write $C^\circ=C\setminus Y$.

Let $V\to X\setminus Y$ be a connected finite \'etale cover
classified by $Q(X,Y)$, and let $W$ be the normalization of $X$
in its function field. It is finite over $X$, normal, integral
and projective. At a point of $C\cap Y$ the strict henselization
of $X$ is regular and the unique boundary divisor is one regular
parameter. Abhyankar's lemma
\cite[XIII, Proposition 5.2, pp.~428--429]{SGA1} trivializes the
cover after adjoining a prime-to-characteristic root of that
parameter. Each local connected component is therefore a cyclic
Kummer chart $z^e=x_1$. Transversality makes its restriction to
$C$ regular of dimension one. Over $C^\circ$ the cover is already
\'etale. Thus the base change $W\times_X C$ is a regular
projective curve.

Apply the connectedness theorem of Martinelli--Naranjo--Pirola
\cite[Theorem 1.1, p.~31]{MNP} to $W\to\mathbf P^r$.
Its image has dimension $d$, while $\dim L=r-d+1$, so the
theorem makes $W\times_X C$ connected. A connected regular
curve is integral; removing the finitely many points above $Y$
leaves a connected tame cover of $C^\circ$. The same curve
therefore tests every connected finite $Q(X,Y)$-cover. It follows
that $\pi_1^t(C^\circ)\to Q(X,Y)$ is surjective. The tame-curve
theorem makes its source finitely generated and proves the lemma.
\end{proof}

\begin{proof}[Proof of Theorem~\ref{thm:boundary}]
Let $N\subset X$ be the prescribed quasi-projective neighborhood
of $Y$. Raynaud's Chow lemma gives a blowup with center outside
$N$ whose source is quasi-projective
\cite[Chapter 4, \S4, Proposition 5, p.~30]{Raynaud}.
Properness makes the source projective. Its normalization is
finite and still projective. We obtain a normal projective
modification $q:X'\to X$ which is an isomorphism over all of
$N$. Put $Y'=q^{-1}(Y)$ and $U'=X'\setminus Y'$.
Every divisorial boundary and its valuation ring are unchanged.

The map $\pi_1(U')\to\pi_1(X\setminus Y)$ is surjective.
Indeed a connected finite \'etale cover of a normal integral
scheme is integral. Its generic fiber is a field, which stays
the same under a birational pullback; the pullback over normal
integral $U'$ is therefore connected. Tameness at the unchanged
boundary valuations gives a surjection
\[
 Q(X',Y')\twoheadrightarrow Q(X,Y).
\]
Lemma~\ref{lem:projective} proves finite generation of the
source and hence of the target. If the boundary has codimension
at least two, there are no divisorial inertia groups to quotient
out, giving the stated assertion about the full fundamental group.
One closed point has an affine neighborhood, so the point case
follows.

For a possibly nonnormal proper variety $Z$, first remove its
nilpotents, which do not change its \'etale site, and take the
disjoint union of the normalizations of its irreducible components.
This is a finite proper surjection. The inverse image of the
specified neighborhood is quasi-projective, and the inverse
image of the omitted set has codimension at least two in each
component. The normal case gives finite generation on each
component of the punctured normalization. Effective proper
descent for finite \'etale covers, in the finite-generation form
of \cite[IX, Corollary 5.2, pp.~247--248]{SGA1}, proves finite
generation downstairs. The finitely many connected components
of the descent relation contribute only finitely many extra
generators.
\end{proof}

\begin{proposition}\label{prop:finite-wild}
Let $X$ be normal proper integral and $Y\subsetneq X$ closed,
over an algebraically closed field. Suppose a normal projective
modification $X'\to X$ has finite wild inertia image in $Q(X,Y)$
at every divisorial component of $X'\setminus q^{-1}(X\setminus Y)$.
Then $Q(X,Y)$ is finitely generated. It is enough to check the
exceptional boundary divisors.
\end{proposition}
\begin{proof}
There are finitely many boundary divisors and finitely many
nonidentity elements in the union of their wild images. A
profinite group has an open normal subgroup $H$ missing all
those elements. Let $V\to X\setminus Y$ be the connected Galois
cover with quotient $Q(X,Y)/H$. Normalize $X'$ in its function
field, obtaining a normal projective variety $T$. Above
$U'=q^{-1}(X\setminus Y)$ its open $V'$ is the connected
birational pullback of $V$. As in the preceding proof,
$\pi_1(V')\to\pi_1(V)\to H$ is surjective.

For compatible prolongations of a discrete valuation under a
finite separable extension $L/K$, absolute inertia and wild
inertia satisfy
\[
 I_w=I_v\cap G_L,\qquad P_w=P_v\cap G_L.
\]
For the second identity the intersection is pro-$p$, while its
quotient in $I_w$ is a subgroup of the prime-to-$p$ tame inertia
quotient. This characterization also applies to imperfect
residue fields \cite[Tag 0BUA, Lemma 58.13.5]{Stacks}.
Every boundary divisor of $T$ lies over a boundary divisor of
$X'$, because $T\to X'$ is finite. Its wild image in $H$ is
therefore trivial, by the choice of $H$ and its normality.
The displayed surjection factors through $Q(T,T\setminus V')$.
Lemma~\ref{lem:projective} makes $H$ finitely generated, and
the finite extension $Q(X,Y)$ is then finitely generated.

The strict transform of an original boundary divisor has the
same valuation ring; its wild image was killed in the definition
of $Q(X,Y)$. Only exceptional boundary divisors need examination.
\end{proof}

\subsection{Comparison with valuation tameness}
Achinger--H\"ubner--Lara--Stix prove that $\pi_1^t(U/k)$ is
topologically finitely generated for every connected scheme $U$
of finite type over an algebraically closed field
\cite[Theorem 2.6.1, p.~23]{AHLS}. Their definition tests
valuations of residue fields over $k$, including valuations whose
centres on a proper compactification have codimension at least two
\cite[Definitions 2.2.1--2.2.2, p.~12]{AHLS}. Thus
\[
 \pi_1(U)\twoheadrightarrow Q(X,Y)
 \twoheadrightarrow\pi_1^t(U/k),\qquad U=X\setminus Y.
\]
Finite generation of the last quotient does not settle the question
for $Q(X,Y)$. The following example makes the difference explicit.

\begin{example}\label{ex:hidden-wild}
Over an algebraically closed field of characteristic $p>0$, there
is a normal projective surface $X$ with a closed point $q$ such that
\[
 Q(X,\{q\})=\pi_1(X\setminus\{q\})\simeq\Z/p\Z,
 \qquad \pi_1^t((X\setminus\{q\})/k)=1.
\]
An exceptional divisorial valuation centred at $q$ has wild inertia
image equal to $\Z/p\Z$.
\end{example}
\begin{proof}
Let $G=(\F_p,+)$ act on $S=\mathbf P^1\times\mathbf P^1$ by
$a(x,y)=(x+a,y+a)$. Every nonidentity element fixes exactly
$P=(\infty,\infty)$, so the action is free on
$S^\circ=S\setminus\{P\}$. The finite quotient $X=S/G$
is normal and projective. If $q$ is the image of $P$, then
$S^\circ\to U=X\setminus\{q\}$ is a connected finite \'etale
$G$-cover. The projective line is simply connected by
Riemann--Hurwitz, and the proper homotopy exact sequence makes
$S$ simply connected \cite[Section 58.15, Tag 0BUM]{Stacks}.
Purity identifies $\pi_1(S^\circ)$ with $\pi_1(S)$
\cite[X, Th\'eor\`eme 3.1]{SGA1}, so the Galois-cover exact
sequence gives $\pi_1(U)=G$. There is no boundary divisor on $X$,
and hence $Q(X,\{q\})=G$.

Blow up $S$ at $P$ and write $E$ for the exceptional divisor.
The action lifts. In the parameters $u=1/x$, $v=1/y$ it sends
$u$ to $u/(1+au)$ and $v$ to $v/(1+av)$. On the blowup chart
$v=ut$ the second coordinate becomes
\[
 t\longmapsto t\frac{1+au}{1+aut}.
\]
Thus $G$ fixes $E$ and acts trivially on its residue field.
For the valuation of $E$ in $k(S)/k(X)$, both decomposition
and inertia are all of $G$; since $G$ is a $p$-group,
wild inertia is also all of $G$. Its restricted valuation is divisorial on the
normal projective modification $\operatorname{Bl}_P(S)/G\to X$
and is centred at $q$. The all-valuation tame quotient must kill
this whole image, so it is trivial. This computation uses inertia,
not just a ramification index, and also allows an inseparable
residue extension.
\end{proof}
